
\subsection{Case Study 3: Fuel Consumption and Emissions}
Main propulsion fuel use was calculated using the regression-based power model. Results showed that the longer Arctic route incurred $\sim$875 tonnes of fuel versus $\sim$756 tonnes for the Suez route, with corresponding \ac{co2} emissions of $\sim$2,724 tonnes and $\sim$2,356 tonnes respectively. This confirms that distance differences propagate directly into sustainability assessments, amplifying the impact of routing assumptions.

\subsection{Case Study 4: Auxiliary Engine Consumption}
Auxiliary loads, assumed at 10\% of installed power, scaled proportionally with voyage time. Results indicated an additional 161 tonnes of fuel and 502 tonnes of \ac{co2} emissions for the Arctic route compared to 139 tonnes and 434 tonnes for the Suez route. Although auxiliary consumption is smaller than main propulsion, its contribution is non-negligible over long voyages.

\subsection{Case Study 5: Sensitivity Analysis}
A parametric sweep over service speed and \ac{dwt} confirmed the expected cubic scaling of power with speed. Fuel and emissions differences between the Arctic and Suez routes widened sharply at higher speeds, while vessel size exerted a more linear effect. These results emphasize that operational choices interact strongly with methodological assumptions.

\subsection{Summary of Findings}
Table~\ref{tab:results_summary} consolidates the key numerical results across case studies, while Figures~\ref{fig:fuel_comparison}--\ref{fig:3d_sensitivity} provide graphical comparisons. The central finding is that routing methodology fundamentally shapes the estimated efficiency of Arctic shipping: while GC routing suggests large savings, sea-only routing demonstrates higher distance, fuel use, and emissions for the Arctic route under identical assumptions.


